% this file is called up by thesis.tex
% content in this file will be fed into the main document

%: ----------------------- introduction file header -----------------------
\chapter{Related Works}
\label{ch:related-works}

% the code below specifies where the figures are stored
\ifpdf
    \graphicspath{{2_related_works/figures/PNG/}{2_related_works/figures/PDF/}{2_related_works/figures/}}
\else
    \graphicspath{{2_related_works/figures/EPS/}{2_related_works/figures/}}
\fi


\section{Wind Turbine Condition Monitoring}
\label{sec:wind-turbine-condition-monitoring}
\gls{wtcm} plays a central role in reducing operational costs and improving the reliability of both onshore and offshore wind farms.
As Dao et al.
(2019) note, operational and maintenance costs account for approximately 20-30\% of the Levelised Cost of Energy, making early fault detection a key strategy for cost control.
Condition monitoring enables operators to identify incipient failures, minimize unplanned downtime, and optimize energy production through targeted maintenance.
\cite{yang_operations_2021} similarly emphasize that proactive monitoring improves asset longevity and reduces corrective maintenance needs, with particular benefits for offshore installations where accessibility is constrained.

Effective \gls{wtcm} integrates multiple data sources, including \gls{scada} systems, \gls{cms}, and historical failure logs.
\gls{scada} data—capturing variables such as power output, rotor speed, vibration, and temperature—facilitates real-time anomaly detection, while \gls{cms} focuses on mechanical components such as gearboxes, bearings, and drive shafts.
The combination of these sources supports not only short-term operational decisions but also long-term reliability modelling.


\subsection{Traditional Statistical Approaches}
\label{subsec:traditional-statistical-approaches}
Early approaches to \gls{wtcm} relied heavily on statistical models for anomaly detection and degradation trend analysis.
ARIMA models were applied to capture temporal dependencies in \gls{scada} data, particularly for forecasting power output or component temperatures under normal operating conditions~\cite{chesterman_overview_2023}.
\gls{pca} was used to reduce the dimensionality of multivariate datasets and to extract latent features associated with turbine health status.
In parallel, threshold-based monitoring systems compared sensor readings against predefined limits, triggering alarms when deviations exceeded acceptable ranges.

While these statistical methods provided valuable insights in the early stages of \gls{wtcm}, they suffer from notable limitations.
Many approaches assume linear relationships between variables, which rarely hold in the complex, nonlinear dynamics of turbine operation.
Moreover, their scalability to modern high-frequency, high-dimensional \gls{scada} datasets is limited—particularly as offshore wind farms generate terabytes of heterogeneous sensor data.
Threshold-based approaches, in particular, are prone to false alarms when environmental variability (e.g., wind speed fluctuations) is not adequately modelled.



\subsection{Transition to Machine Learning Techniques}
\label{subsec:transition-to-machine-learning-techniques}
The statistical foundations of \gls{wtcm} have directly informed the transition to \gls{ml}-based approaches.

For example, \gls{pca}'s dimensionality reduction principles underpin the structure of autoencoders used for unsupervised anomaly detection.
Similarly, ARIMA's time series modelling concepts inspire \glspl{rnn} and temporal convolutional networks that capture complex, nonlinear dependencies.
Threshold-based methods have evolved into supervised learning classifiers that learn optimal decision boundaries from historical fault data, rather than relying on fixed limits.
As turbine datasets have grown in size and complexity, these \gls{ml} approaches have demonstrated improved fault detection accuracy, adaptability to nonstationary conditions, and robustness to noisy inputs — making them a natural evolution of earlier statistical techniques.


\subsection{Conclusion}
Statistical methods laid the foundation for modern \gls{wtcm} by providing initial frameworks for anomaly detection and fault diagnosis.
However, their inherent limitations in handling non-linearity and high-dimensional data have necessitated the shift toward \textbf{data-driven machine learning approaches}.
The integration of \gls{ml} techniques, such as autoencoders and \glspl{lstm}, has significantly enhanced the accuracy and scalability of wind turbine monitoring, enabling more reliable and cost-effective maintenance strategies.
Future research should focus on hybrid models that combine statistical robustness with \gls{ml} adaptability to further optimize wind turbine performance and reduce maintenance costs.


\section{Machine Learning in Normal Behavior Models}
\label{sec:machine-learning-in-normal-behavior-models}